%%% SECTION 2: U.S. GOVERNMENT FRAMEWORK FOR PHYSICAL AI TRIALS %%%

Physical AI oncology trials must be understood within the full constitutional
and statutory framework of U.S. government, spanning the Legislative, Executive,
and Judicial branches. The governance structure for Physical AI builds on
existing federal branch operations rather than requiring new institutional
frameworks, leveraging established institutional credibility and legal authority.
This section adapts the comprehensive research on U.S. federal branch operations
statutes and their application to AI-enabled clinical trials to Physical AI
oncology trial contexts.

\subsection{Constitutional Foundation}
\label{subsec:constitutional}

The U.S. federal government's three-branch structure is established primarily by
the U.S. Constitution. Article I vests legislative powers in a Congress
consisting of a Senate and House of Representatives. Article II vests executive
power in a President, with the Appointments Clause establishing Senate
advice-and-consent for officers and the Take Care duty (Article II, Section 3)
providing the constitutional anchor for executive implementation and supervision
of federal law. Article III vests the judicial power in one Supreme Court and
such inferior courts as Congress may establish.

For Physical AI oncology trials, this three-branch structure provides essential
checks and balances. Legislative authority creates the statutory framework
through which Physical AI trial programs are funded and overseen. Executive
authority implements that framework through FDA and HHS, which directly
regulate clinical trial conduct. Judicial authority resolves disputes and
interprets the boundaries of regulatory action. The Congressional Review Act
(5 U.S.C. \S~801 et seq.) provides legislative oversight of executive
rulemaking, meaning that any FDA regulations specific to Physical AI would be
subject to congressional review. The Administrative Procedure Act (5 U.S.C.
\S~551 et seq.) requires that Physical AI regulations follow standard
notice-and-comment rulemaking procedures, ensuring public participation and
transparency.

The Appointments Clause has practical implications for Physical AI governance:
FDA commissioners and senior officials who would oversee Physical AI trial
approvals are appointed through the constitutionally mandated process, ensuring
democratic accountability for regulatory decisions affecting patient safety.
The Take Care duty requires the executive branch to faithfully implement
clinical trial statutes, preventing either over-regulation that would block
beneficial Physical AI innovations or under-regulation that would compromise
patient safety.

\subsection{Legislative Branch and Physical AI Regulation}
\label{subsec:legislative}

The legislative branch's role in Physical AI regulation operates through several
operational statutes. The Congressional Budget and Impoundment Control Act of
1974 (codified across titles, with CBO at 2 U.S.C. \S~601) creates the
Congressional Budget Office and statutory budget procedures that structure
annual budgeting and congressional scorekeeping. FDA funding for Physical AI
trial review capacity depends directly on congressional appropriations authority
exercised through these budget procedures.

The Congressional Review Act (5 U.S.C. \S~801 et seq.) establishes a mechanism
for Congress to review and disapprove certain agency rules under defined
procedures. If FDA were to promulgate rules specifically governing Physical AI
clinical trials, these rules would be subject to congressional review,
providing a check on executive rulemaking power. The Congressional
Accountability Act of 1995 (2 U.S.C. ch. 24, \S~1301; Pub. L. 104-1) applies
workplace and labor protections to legislative-branch employing offices,
establishing precedent for how new categories of workers (including those
overseeing Physical AI systems) receive legislative protection.

The Government Accountability Office (GAO), with Comptroller General
authorities established under 31 U.S.C. \S~702, provides audit and
investigation capacity for oversight of Physical AI trial programs. GAO audits
of FDA's review processes for AI-enabled medical products would apply equally
to Physical AI trial oversight, ensuring that congressional appropriations are
used effectively and that regulatory processes meet statutory requirements.

Congressional appropriations directly affect FDA's capacity to review and
approve Physical AI trial applications. If Physical AI trials require new
categories of expertise (robotics engineers, AI safety specialists,
cybersecurity analysts) within FDA review teams, congressional funding must
support these positions. The budget process statutes ensure that these resource
needs are evaluated through established scorekeeping procedures rather than
ad hoc allocation.

\begin{table}[H]
\centering
\caption{Core Legislative Statutes Applicable to Physical AI Trial Governance}
\label{tab:legislative-statutes}
\small
\begin{tabularx}{\textwidth}{l X l}
\toprule
\textbf{Statute} & \textbf{Physical AI Relevance} & \textbf{Citation} \\
\midrule
Budget and Impoundment Control Act & FDA funding for Physical AI review capacity & 2 U.S.C. \S~601 \\
Congressional Review Act & Review of Physical AI-specific FDA rules & 5 U.S.C. \S~801 \\
Congressional Accountability Act & Worker protections for Physical AI oversight staff & 2 U.S.C. \S~1301 \\
GAO/Comptroller General & Audit of Physical AI trial programs & 31 U.S.C. \S~702 \\
Reapportionment Act & House composition for Physical AI policy votes & 2 U.S.C. \S~2a \\
\bottomrule
\end{tabularx}
\end{table}

\subsection{Executive Branch Implementation}
\label{subsec:executive}

The executive branch implements clinical trial regulation through FDA and HHS,
with several foundational statutes governing the process. The Administrative
Procedure Act (APA, Pub. L. 79-404; 5 U.S.C. \S~551 et seq.) defines
``agency'' and related terms, forming the baseline procedural code for federal
administrative action and review. FDA's authority to regulate Physical AI trials
derives from the APA's rulemaking framework: any Physical AI-specific
regulations would follow notice-and-comment procedures, with judicial review
available for arbitrary or capricious agency action.

The Freedom of Information Act (FOIA, Pub. L. 89-487; 5 U.S.C. \S~552)
creates enforceable disclosure duties subject to enumerated exemptions. For
Physical AI trials, FOIA ensures that FDA review processes, approval decisions,
and safety data are accessible to the public, researchers, and industry
participants. The transparency provided by FOIA supports public trust in
Physical AI adoption by making regulatory decision-making visible and
accountable.

The Privacy Act of 1974 (Pub. L. 93-579; 5 U.S.C. \S~552a) establishes
restrictions on disclosure and requirements for systems of records maintained by
federal agencies. When FDA maintains records about Physical AI trial sponsors,
investigators, or participants, the Privacy Act governs how those records may
be used and disclosed. This statute includes criminal penalties for certain
willful violations, providing strong incentives for proper data handling.

The Federal Register Act (44 U.S.C. \S~1505) requires publication of
presidential proclamations and executive orders with general applicability and
legal effect. Any executive orders directing FDA to prioritize Physical AI
trial review or establishing Physical AI governance frameworks must be
published in the Federal Register to have legal effect. The Paperwork Reduction
Act (44 U.S.C. \S~3501) requires OMB review of information collections,
meaning that new data collection requirements for Physical AI trial sponsors
must go through OMB clearance.

The Inspectors General framework (5 U.S.C. ch. 4) requires IG reporting
structures within executive agencies, providing internal oversight and audit
capability for Physical AI trial programs. The Federal Vacancies Reform Act of
1998 (5 U.S.C. \S~3345) establishes rules for acting officials, ensuring
continuity of FDA leadership during transitions that might affect Physical AI
regulatory decisions.

\begin{table}[H]
\centering
\caption{Core Executive Statutes Applicable to Physical AI Trial Operations}
\label{tab:executive-statutes}
\small
\begin{tabularx}{\textwidth}{l X l}
\toprule
\textbf{Statute} & \textbf{Physical AI Relevance} & \textbf{Citation} \\
\midrule
Administrative Procedure Act & Rulemaking for Physical AI regulations & 5 U.S.C. \S~551 \\
Freedom of Information Act & Transparency of Physical AI trial decisions & 5 U.S.C. \S~552 \\
Privacy Act of 1974 & Protection of Physical AI trial records & 5 U.S.C. \S~552a \\
Federal Register Act & Publication of Physical AI executive orders & 44 U.S.C. \S~1505 \\
Paperwork Reduction Act & OMB review of Physical AI data collection & 44 U.S.C. \S~3501 \\
Inspectors General & Internal audit of Physical AI programs & 5 U.S.C. ch. 4 \\
Federal Vacancies Reform & Continuity of FDA Physical AI leadership & 5 U.S.C. \S~3345 \\
Presidential Records Act & Preservation of Physical AI policy records & 44 U.S.C. \S~2201 \\
\bottomrule
\end{tabularx}
\end{table}

The three adapted regulatory standards in this platform demonstrate how
existing executive branch regulatory frameworks can be adapted rather than
replaced. The Adaption: ICH Harmonised Guideline \cite{ich-adapt} builds on the
internationally recognized GCP standard. The Physical AI Adaption: 21 CFR Part
50 \cite{cfr50-adapt} extends existing human subjects protection regulation. The
Physical AI Adaption: 21 CFR Part 312 \cite{cfr312-adapt} extends the IND
application regulation. Each adaptation operates within the APA's rulemaking
framework, demonstrating that Physical AI can be regulated through existing
administrative procedures.

\subsection{Judicial Branch Oversight}
\label{subsec:judicial}

The judicial branch provides the framework for resolving Physical AI regulatory
disputes. The modern statutory backbone for court organization, jurisdiction,
and judicial administration is established in Title 28 of the U.S. Code. The
Supreme Court size and quorum are set by 28 U.S.C. \S~1, with circuit
composition defined by 28 U.S.C. \S~41 (currently thirteen circuits including
the Federal Circuit). The Judicial Conference (28 U.S.C. \S~331) serves as the
judiciary's policy-making body, coordinating administration across all federal
courts.

The Rules Enabling Act (28 U.S.C. \S~2072) authorizes the Supreme Court to
prescribe general rules of practice, procedure, and evidence, while prohibiting
rules that abridge, enlarge, or modify substantive rights. These procedural
rules govern how Physical AI regulatory disputes would be litigated, from
initial filing through appeal. The judicial discipline and complaints process
(28 U.S.C. \S~351) ensures accountability for judges who handle Physical AI
cases.

Courts would review FDA decisions regarding Physical AI trial approvals, safety
actions, and clinical holds under existing administrative law frameworks. The
APA provides for judicial review of agency action, with courts evaluating
whether FDA's Physical AI decisions are arbitrary, capricious, an abuse of
discretion, or otherwise not in accordance with law. This standard of review
means that FDA must articulate reasoned explanations for Physical AI regulatory
decisions, which in turn requires the kind of comprehensive evidence base that
this National Platform provides.

If FDA were to place a clinical hold on a Physical AI trial, the sponsor could
challenge that decision through administrative channels and ultimately through
judicial review. If FDA were to deny an IND application that included Physical
AI components, the denial would be subject to review under the same standards
that apply to all IND decisions. The existing judicial framework thus provides
a safety valve against both over-regulation and under-regulation of Physical AI
trials, ensuring that regulatory decisions are grounded in evidence rather than
speculation.

\subsection{Cross-Branch Governance for Physical AI}
\label{subsec:cross-branch}

Several statutes operate across all three branches to ensure accountability in
Physical AI trial governance. Transparency and recordkeeping requirements
include Federal Register publication requirements (44 U.S.C. \S~1505), FOIA
public access provisions (5 U.S.C. \S~552), agency record preservation
obligations (44 U.S.C. \S~3101), and presidential record definitions
(44 U.S.C. \S~2201). Together, these ensure that Physical AI policy decisions
are documented, preserved, and accessible for public scrutiny.

Budget execution constraints apply through the Antideficiency Act (31 U.S.C.
\S~1341), which prohibits obligations and expenditures exceeding available
appropriations. This means that Physical AI trial programs cannot spend beyond
congressional authorization, providing fiscal discipline for potentially
expensive national deployment efforts. The Ethics in Government Act framework
establishes disclosure and ethics requirements spanning all branches, ensuring
that officials involved in Physical AI regulatory decisions are subject to
conflict-of-interest protections and financial disclosure requirements.

Physical AI oncology trials benefit from the full weight of these existing
cross-branch governance mechanisms. Rather than creating new governance
structures, this platform adapts proven mechanisms to ensure transparency,
fiscal responsibility, and ethical conduct. This approach is consistent with the
platform's broader philosophy of building on established frameworks rather than
replacing them.

\subsection{Current Regulatory Framework for AI in Clinical Trials}
\label{subsec:current-ai-framework}

AI use in clinical trials is regulated primarily through what the AI does rather
than through a dedicated ``AI clinical trials'' statute. FDA guidance
distinguishes four categories of AI use in trial contexts. First, AI used to
generate information or data intended to support regulatory decision-making for
drugs and biologics falls under FDA's January 2025 draft guidance, which
proposes a risk-based credibility assessment framework with seven steps: define
question of interest, define context of use, assess model risk, plan, execute,
document, and determine adequacy. Second, AI embedded in or functioning as
medical device software falls under device lifecycle management guidance and
PCCP requirements. Third, AI used as part of trial conduct and data systems
must comply with 21 CFR Part 11 electronic records requirements and related
computerized systems guidance. Fourth, AI as human subjects research and data
risk falls under Common Rule and IRB review requirements.

\begin{table}[H]
\centering
\caption{Federal Regulations and Guidances Relevant to Physical AI Trials}
\label{tab:regulations-ai}
\small
\begin{tabularx}{\textwidth}{l l X}
\toprule
\textbf{Type} & \textbf{Citation} & \textbf{Physical AI Relevance} \\
\midrule
Regulation & 21 CFR Part 312 & IND for trials with Physical AI drug administration \\
Regulation & 21 CFR Part 50 & Informed consent for Physical AI robotic procedures \\
Regulation & 21 CFR Part 56 & IRB review of Physical AI system risks \\
Regulation & 21 CFR Part 812 & IDE for Physical AI device investigations \\
Regulation & 21 CFR Part 11 & Electronic records from Physical AI systems \\
Regulation & 42 CFR Part 11 & Trial registration including Physical AI methods \\
Regulation & 45 CFR Part 46 & Common Rule for Physical AI human subjects research \\
Regulation & 45 CFR Parts 160, 164 & HIPAA for Physical AI health data \\
Guidance & AI Drug/Biologic (Jan 2025) & AI credibility framework for Physical AI evidence \\
Guidance & AI Device Lifecycle (Jan 2025) & Lifecycle management for Physical AI software \\
Guidance & Decentralized Trials (Sep 2024) & Remote Physical AI monitoring elements \\
Guidance & PCCP for AI Devices (Aug 2025) & Change control for Physical AI system updates \\
Guidance & CDS Software (Jan 2026) & Classification of Physical AI decision support \\
\bottomrule
\end{tabularx}
\end{table}

Physical AI introduces additional considerations across all four categories.
When Physical AI systems generate evidence for regulatory decisions (such as
robotic surgical outcomes data), the credibility assessment must account for
robot-specific variables including mechanical precision, sensor accuracy, and
autonomous decision quality. When Physical AI functions as a regulated device,
the lifecycle management must address hardware degradation, calibration drift,
and physical wear in addition to software updates. When Physical AI generates
electronic trial records (robot telemetry, sensor data, autonomous decision
logs), these records represent new categories that must comply with Part 11
requirements. When Physical AI interacts directly with human subjects, the risk
profile differs fundamentally from software-only AI because of physical contact
and proximity.

\subsection{AI and Oncology: A High-Frequency Use Case}
\label{subsec:ai-oncology}

The FDA Oncology Center of Excellence \cite{fda-oncology-guidance} describes an
Oncology AI Program begun in 2023, linking AI-related FDA guidances used in
oncology regulatory contexts. In oncology trials specifically, AI is frequently
implicated because endpoints and eligibility often rely on imaging (CT, MRI,
PET), digital pathology, complex biomarker patterns, and longitudinal
EHR-derived outcomes, all of which are common AI input modalities.

Physical AI extends these existing AI applications by adding robotic
embodiment. Rather than AI analyzing images to inform human decisions, Physical
AI systems can directly act on imaging analysis by guiding surgical robots to
tumor sites, positioning radiotherapy beams with sub-millimeter precision, or
directing needle-placement robots to biopsy targets. A Cancer Patient's
Journey, Physical AI \cite{patient-journey-paper} documents how this integrated
approach works across all ten stages of an oncology trial, from prescreening
through closeout. The ten-stage pipeline demonstrates that Physical AI is not
limited to isolated procedures but can manage the complete patient experience.

\subsection{Practical Compliance Synthesis}
\label{subsec:compliance-synthesis}

In oncology trials using Physical AI, compliance requires building a defensible
chain from context of use through validation, auditability, and human subjects
protection.

If Physical AI outputs will be used as evidence for FDA decisions about safety,
effectiveness, or quality, sponsors should document the model's context of use,
risk, and credibility assessment consistent with FDA's draft framework. The
adapted 21 CFR Part 312 \cite{cfr312-adapt} provides the IND framework for
including Physical AI evidence in regulatory submissions, while the adapted ICH
E6(R3) \cite{ich-adapt} establishes the quality management system for ensuring
Physical AI data integrity.

If Physical AI functions as a regulated device, the trial may implicate IDE and
device marketing pathways, with lifecycle management and PCCP expectations
governing what changes are allowed and how they are controlled. The USL scoring
framework \cite{usl-standard} provides quantitative assessment of each robot's
readiness for clinical deployment, with Dimension D (clinical trial
collaboration) directly measuring regulatory compliance status.

If Physical AI relies on decentralized or digital health technology data, the
data attribution, metadata integrity, and Part 11 boundaries become critical
compliance points. The MCP server infrastructure \cite{national-mcp-paper}
addresses this through standardized data flow pathways and conformance level
requirements that ensure data integrity across distributed Physical AI systems.

If Physical AI training and validation rely on secondary data and research
exemptions, IRB considerations around identifiability and bias must be
addressed. The federated learning framework \cite{fl-paper} provides
privacy-preserving approaches to model training that reduce identifiability
risks while enabling multi-site data utilization under HIPAA \cite{hipaa} and
state privacy law requirements.

The adapted 21 CFR Part 50 \cite{cfr50-adapt} ensures that human subjects
protection extends to all Physical AI interactions, with informed consent
requirements covering robot descriptions, AI decision-making transparency,
patient rights to refuse robotic procedures, and emergency stop procedures.
Together, the three adapted standards provide comprehensive responses to each
compliance pathway, demonstrating that Physical AI can operate within the
existing regulatory architecture with targeted adaptations.
